\chapter*{\center{Introduction Générale}}
\addcontentsline{toc}{chapter}{Introduction Générale}

Avec l'évolution rapide de l'intelligence artificielle, le secteur de l'éducation vit aujourd'hui une transformation sans précédent. Parmi les innovations marquantes, la convergence entre la vision par ordinateur et les modèles de langage de grande taille (\textit{LLMs}) se démarque par sa capacité à simuler un accompagnement humain tout en garantissant une rigueur académique infaillible. Ces systèmes ne sont plus de simples outils techniques, mais des piliers centraux de l'accompagnement pédagogique, notamment sur les plateformes de certification en ligne.

Le projet que je présente, intitulé \textbf{Certif-fun}, s'inscrit dans cette dynamique d'innovation globale. Il vise à concevoir un écosystème intelligent dédié non seulement à la \textbf{gestion des certifications} et à la \textbf{surveillance automatisée}, mais aussi à l' \textbf{encadrement structuré} et aux \textbf{réunions virtuelles} à distance. À travers une combinaison harmonieuse entre le \textit{proctoring} par IA et la \textbf{génération de contenu pédagogique}, Certif-fun tente de combler le fossé entre la surveillance sécurisée et le soutien collaboratif.

L'objectif principal de ce travail est donc multidimensionnel : d'une part, développer un module de supervision proactive s'appuyant sur YOLOv8 et MediaPipe pour identifier les comportements suspects ; d'autre part, exploiter la puissance d'Ollama pour générer instantanément des quiz et des supports de cours. Ce projet vise également à unifier l'espace de certification avec un module de visioconférence et d'encadrement, améliorant ainsi l'expérience globale des formateurs et des apprenants.

Pour atteindre ces buts, une méthodologie rigoureuse a été adoptée. Elle débute par une exploration approfondie de l'état de l'art, analysant les recherches antérieures menées sur les plateformes éducatives existantes ainsi que les avancées récentes en IA générative et en vision par ordinateur. Ensuite, une phase de modélisation a permis de structurer cet écosystème à l'aide des diagrammes UML et BPMN. Enfin, le développement s'est appuyé sur une architecture microservices performante (Spring Boot, React, FastAPI), garantissant la scalabilité du système.

À travers ce travail, mon objectif est de mettre en œuvre une preuve de concept robuste démontrant la faisabilité technique et pédagogique d'un tel écosystème. Le développement réalisé aspire à poser les fondations d'un système évolutif, capable d'unifier la sécurité des examens et le soutien pédagogique au sein d'une architecture moderne, ouvrant ainsi de nouvelles perspectives pour l'avenir des plateformes éducatives numériques.
